\documentclass[12pt]{article}

\usepackage[a4paper,top=2cm,bottom=2cm,left=2cm,right=2cm,marginparwidth=1.5cm]{geometry}
\usepackage{graphicx}
\usepackage{amssymb}
\usepackage{epstopdf}
\usepackage[german]{babel}
\usepackage[utf8]{inputenc}
\usepackage{pdfpages}
\usepackage{fancyhdr}
\DeclareGraphicsRule{.tif}{png}{.png}{`convert #1 `dirname #1`/`basename #1 .tif`.png}

\title{Bericht}
\author{Alisa Pesteritz, Saliha Altan, Wiktoria Cholewa, \\ 
	Dominik Falk, Andreas Kasarinow, Johanna Kleidorfer, \\
	Adrika Narasimhan, Rosina Röckseisen}

\begin{document}

\newpage
\maketitle

\newpage
\pagestyle{fancy}
\fancyhead{}
\fancyhead[L]{Alle}
\section{Einleitung}
\label{einleitung}

\section{Arbeitseinteilung (oder sowas)}
\label{arbeitseinteilung}
% TODO: wer hat was gemacht und warum?
% Datenfluss (von Johanna) einfügen? (oder zumindest beschreiben?)

% ----------------------------------------------------------------------------------------------
\newpage
\pagestyle{fancy}
\fancyhead{}
\fancyhead[L]{Alisa}
\section{Pathways und Dateneinlesen}
\label{pathways}

% ----------------------------------------------------------------------------------------------
\newpage
\pagestyle{fancy}
\fancyhead{}
\fancyhead[L]{Rosina \& Wiki}
\section{Normalisierung}
\label{normalisierung}


\section{Clusteranalyse}
\label{clusteranalyse}
% TODO: theoretische Erklärung des Algorithmus

\subsection{Single Linkage}
\label{single-linkage}

\subsection{Average Linkage}
\label{average-linkage}

\subsection{Complete Linkage}
\label{complete-linkage}

% ----------------------------------------------------------------------------------------------
\newpage
\pagestyle{fancy}
\fancyhead{}
\fancyhead[L]{Saliha \& Domi}
\section{Visualisierung}
\label{visualisierung}

\subsection{Heatmap}
\label{heatmap}

\subsection{Dendrogramm}
\label{dendrogramm}

% ----------------------------------------------------------------------------------------------
\newpage
\pagestyle{fancy}
\fancyhead{}
\fancyhead[L]{Adrika \& Andreas}
\section{GUI}
\label{gui}

\subsection{Serverlogik}
\label{serverlogik}

\begin{table}[]
	\begin{tabular}{l|l|l}
		Aufgabe                       & Zuständig & Arbeitsaufwand \\ 
		\hline
		Vorbereitung des Datensatzes  & Andreas   & 10 Stunden     \\
		Preset-Funktion               & Andreas   & 5 Stunden      \\ 
		Cache                         & Andreas   & 7 Stunden      \\ 
		PDF-Export			          & Andreas   & 10 Stunden     \\
		Besprechungen 				  & Andreas   & 20 Stunden     \\
	\end{tabular}
\end{table}

\subsubsection{Verarbeitung des Datensatzes}
\label{Verarbeitung des Datensatzes}
Da nicht davon ausgegangen werden konnte, dass alle hochgeladenen Datensätze vollständig und ohne Fehler hochgeladen werden, wurde die Dateneingabe um eine automatische NA-Prüfung erweitert mit der fehlende Daten erkannt werden und die verhindert, dass fehlerhafte Werte in die Clusteranalyse übergeben werden mit denen sie nicht umgehen kann. Nach dem Upload wird berechnet, wie viele NA-Werte vorhanden sind und in welchen Zeilen und Spalten sie auftreten. Zeilen und Spalten, die vollständig aus fehlenden Werten bestehen, werden automatisch entfernt. Bei Zeilen oder Spalten mit mindestens 50 Prozent fehlenden Werten wird dagegen eine Entscheidung des Nutzers verlangt. Dabei können die betroffenen Zeilen oder Spalten entfernt oder beibehalten werden. Verbleibende NA-Werte werden anschließend durch Zeilenmittelwerte ersetzt. Die Informationen über die entfernte Zeilen, entfernte Spalten und ersetzte Werte werden in der GUI angezeigt. Die Implementierung dieser Funktionalität war eher schwierig, weil sich die Anforderungen während der Entwicklung mehrfach änderten bzw. man festgestellt hat das bestimmte Sachen nicht bedacht wurden oder auch schlecht kommuniziert wurde. Zu Beginn wurde die Möglichkeit der manuellen Behandlung jedes Wertes eingebaut, die in der Theorie sinnvoll erschien, später zeigte sich jedoch, dass die Lösung die Anwendung hauptsächlich komplizierter machten, ohne einen deutlichen Mehrwert zu bieten. Teilweise wurde viel Zeit in Ideen investiert, die später wieder vereinfacht oder vollständig entfernt wurden. Das automatische Entfernen einer Zeile mit vielen fehlenden Werten ist beispielsweise einfach umzusetzen, kann jedoch zum Verlust relevanter Gene führen. Daher wurde für kritische Fälle, oberhalb der 50 Prozent, eine Nutzerentscheidung eingeführt
\subsubsection{Preset-Funktion}
\label{Preset-Funktion}
Die Preset-Funktion wurde implementiert, um Einstellungen dauerhaft speichern und später einfacher wiederverwenden zu können. Ohne diese Funktion müssten die verschiedenen Parameter wie Clusterverfahren, die Normalisierung, die Distanzmethode, die Farbpalette, verschiedene numerische Parameter und die ausgewählten Pathways bei jeder Sitzung erneut eingestellt werden. Dies wäre besonders bei größeren Konfigurationen sehr zeitaufwendig und fehleranfällig. Die Einstellungen werden zunächst in einer Liste zusammengefasst und mithilfe des Pakets \texttt{jsonlite} als JSON-Datei gespeichert. Beim Laden wird die Datei wieder in eine Liste umgewandelt. JSON wurde gewählt, da das Format leicht gespeichert, eingelesen werden kann. Ein Preset enthält nur die Einstellungen der Analyse. Die größte Schwierigkeit bestand nicht im Speichern der JSON-Datei, sondern beim Laden des Presets. Es gibt viele Eingaben die konditional sind wie Minkowski mit dem P-Wert oder die Parameter alpha, beta und gamma ausschließlich für das  Linkage-Verfahren benötigt werdender dann hinzugefügt werden mussten. Einige Einstellungen sind außerdem mehrfach in der Benutzeroberfläche vorhanden, beispielsweise auf der Parameterseite und zusätzlich in der Sidebar. Beim Laden eines Presets mussten deshalb nicht nur der Zustand der "Parameter wählen" Seite geupdatet werden sonder auch die zugehörigen Eingabefelder beim "Visualisierung" Reiter upgedatet werden, sonst hätte die GUI andere Werte anzeigt als für die Berechnung verwendet worden wären. Die Pathways stellten einen weiteren Sonderfall dar, da die verfügbaren Werte zunächst aus der Datenbank geladen werden müssen.

\subsubsection{Cache}
\label{Cache}
Die Berechnung der Distanzmatrizen ist insbesondere bei größeren Datensätzen rechen- und zeitintensiv. Da sowohl eine Distanzmatrix für die Patienten als auch eine Distanzmatrix für die Gene benötigt wird, kann dieser Schritt einen großen Teil der Laufzeit ausmachen. Der Cache wurde deshalb implementiert, um bereits berechnete Matrizen wiederzuverwenden, wenn sich weder die Daten noch die entscheidenden Einstellungen geändert haben. Der Cache speichert die beiden Distanzmatrizen sowie einen Schlüssel, der die aktuelle Berechnungseinstellungen beinhaltet. Dieser Schlüssel berücksichtigt unter anderem die Datenmatrix, die Normalisierung, die Distanzmethode, die Pathway-Auswahl und wenn ausgewählt, den Minkowski-Parameter. Vor einer neuen Berechnung wird geprüft, ob der aktuelle Schlüssel mit dem gespeicherten Schlüssel übereinstimmt. Ist dies der Fall, können die vorhandenen Matrizen verwendet werden. Andernfalls werden sie neu berechnet und anschließend im Cache gespeichert. Die Schwierigkeit bestand hier darin zu entscheiden wann der Cache geleert werden muss und wann dieser noch gültig ist. Wird dies falsch entschieden können alte Distanzmatrizen mit neuen Daten oder Einstellungen kombiniert werden und dadurch ein fehlerhafter Report entsteht. Außerdem müssen hier wieder alle Parameter die in der GUI geändert werden können aufgenommen werden, da es hier wieder zwei Möglichkeiten gibt, die Parameter wählen-Seite und die Heatmap-Seite. Daher musste für jede Einstellung unterschieden werden, ob sie die Berechnung oder lediglich die Visualisierung beeinflusst oder die grundlegende Berechnung. Eine Änderung der Farbpalette benötigt beispielsweise keine neue Distanzmatrix. Eine andere Normalisierung, Distanzmethode oder Pathway-Auswahl verändert diese, auch der Datensatz selbst muss geprüft werden.

\subsubsection{PDF-Export}
\label{PDF-Export}

Der PDF-Export wurde implementiert, um die Clusterergebnisse außerhalb der Anwendung anschauen und weitergeben zu können. Der PDF-Report enthält die verwendeten Parameter sowie das Patienten-Dendrogramm, das Gen-Dendrogramm und die Heatmap. Zu Beginn wurde versucht, den Bericht vollständig mit dem Paket \texttt{tinytex} zu erzeugen. Dieses funktionierte leider nicht und nach mehrmaligen versuch gab es immer einen Fehler aus. Anstatt die Ursache weiter zu suchen, wurde nach einer eigenen Lösung gesucht. Dabei werden die Parameterseite und die einzelnen Visualisierungen getrennt als PDF-Dateien erzeugt. Danach werden die Dateien sortiert und mit dem Paket \texttt{qpdf} zu einem gemeinsamen Bericht zusammengefügt. Dieser Ansatz war wesentlich aufwendiger als ursprünglich geplant, da die Parameterseite mit dem Base-R-Grafiksystem manuell gestaltet werden musste. Gleichzeitig mussten unterschiedliche Seitengrößen berücksichtigt werden. Besonders die Heatmap stellte bei großen Datensätzen mit mehreren hundert langen Gennamen ein Problem dar. Dafür wurden sowohl mehrseitige Darstellungen als auch eine einzelne große PDF-Seite getestet, auch die Reihenfolge musste stimmen bei der  Zusammenkunft. Auch die Übertragung der Farbpalette und der korrekten Gennamen erforderte weitere Anpassungen. Rückblickend wäre es vermutlich einfacher gewesen, das Paket \texttt{tinytex} zum Laufen zu bringen und dessen Fehler zu beheben.
\newpage
\pagestyle{fancy}
\fancyhead{}
\fancyhead[L]{Johanna}
\subsection{Farben}


\end{document}
